\documentclass[aspectratio=169, 10pt]{beamer}
\usepackage{beamer_theme_cienciadados2026}

\title[Aula 13 - Seaborn Estatístico]{Ciência dos Dados I: Análise Exploratória}
\subtitle{Aula 13: Visualização de Dados II: Análise Estatística com Seaborn}
\author{Prof. Dr. Ney Lemke}
\institute{UNESP -- Instituto de Biociências de Botucatu}
\date{Versão 2026}

\begin{document}

\frame{\titlepage}

\begin{frame}{Sumário da Aula}
  \tableofcontents
\end{frame}

\section{Filosofia do Seaborn}

\begin{frame}{Visualização Estatística de Alto Nível}
  Construído sobre o Matplotlib com integração nativa às estruturas do Pandas:
  \begin{itemize}
    \item \textbf{Mapeamento Semântico}: Mapeia colunas do DataFrame diretamente para atributos visuais (cor com \texttt{hue}, tamanho com \texttt{size}, formato com \texttt{style}).
    \item \textbf{Estatística Embutida}: Estima automaticamente intervalos de confiança, ajusta distribuições e calcula agregações.
    \item \textbf{Temas Elegantes}: Paletas cromáticas perceptualmente uniformes e legíveis.
  \end{itemize}
\end{frame}

\section{Distribuições: Histogramas, KDE e Boxplots}

\begin{frame}[fragile]{Análise Univariada e Detecção de Outliers}
  \begin{columns}
    \begin{column}{0.5\textwidth}
      \begin{block}{Histograma e KDE}
        Permite enxergar bimodalidades, assimetrias (skewness) e caudas pesadas:
      \end{block}
      \begin{lstlisting}
sns.histplot(df['price'], kde=True)
      \end{lstlisting}
    \end{column}
    \begin{column}{0.48\textwidth}
      \begin{block}{Boxplot / Violin Plot}
        Exibe a mediana, intervalo interquartil (IQR) e identifica outliers:
      \end{block}
      \begin{lstlisting}
sns.boxplot(x='body-style', y='price', data=df)
      \end{lstlisting}
    \end{column}
  \end{columns}
\end{frame}

\section{Relações Multivariadas e Mapas de Calor}

\begin{frame}[fragile]{Matrizes de Correlação e Heatmaps}
  Visualizando o grau de associação linear entre múltiplos atributos quantitativos simultaneamente:
  \begin{lstlisting}
corr = df[['price', 'horsepower', 'city-mpg', 'engine-size']].corr()

sns.heatmap(corr, annot=True, cmap='coolwarm', fmt='.2f', vmin=-1, vmax=1)
  \end{lstlisting}
  \begin{block}{Interpretação Exploratória}
    - Correlação positiva forte ($\approx +0.8$): à medida que a potência cresce, o preço também sobe. \\
    - Correlação negativa forte ($\approx -0.7$): veículos mais pesados e potentes consomem mais combustível (menor MPG).
  \end{block}
\end{frame}

\begin{frame}{Resumo da Aula}
  \begin{itemize}
    \item Seaborn reduz dezenas de linhas de código Matplotlib a chamadas semânticas expressivas.
    \item Boxplots e heatmaps são ferramentas indispensáveis de triagem exploratória.
    \item \textbf{Prática}: Distribuições estatísticas, cruzamentos multivariados e matrizes de correlação.
  \end{itemize}
\end{frame}

\end{document}
